% ============================================
% DOC-26
% Extraído de main.pandoc.tex
% ============================================
\section*{\textbf{DOC-26} --- Anotações de Oliveira — O Mecanismo}
\markboth{DOC-26 --- Anotações de Oliveira — O Mecanismo}{}

\noindent \textit{Documento extraído do arquivo pericial.}

\vspace{0.5em}\noindent\rule{\textwidth}{0.5pt}\vspace{0.5em}

\noindent {\footnotesize\texttt{ARQUIVO P\'UBLICO $\cdot$ PROVA DOCUMENTAL $\cdot$ PROCESSO 1.260}}\\[-0.2em]

\noindent\textbf{Documento:} folhas avulsas encontradas no bolso interno do caderno pessoal do Dr. Heitor Oliveira, amarradas com barbante. \textbf{Data estimada:} junho de 1979 a março de 2026. \textbf{Estado:} as folhas mais recentes estão escritas em papel diferente, com caligrafia trêmula e envelhecida --- como se tivessem sido escritas décadas depois, mas amarradas ao mesmo maço.

\subsubsection*{Folha 1 --- 19 de junho de 1979}

\noindent Descobri hoje o que o diário faz com o tempo.

\noindent Não é só escrever. É \textit{endereçar\textit{. Quando escrevo no diário de Joaquim, posso escolher }para quando} aquela palavra vai chegar. Não é controle consciente --- é como entoar uma nota certa e ver o vidro vibrar. A página vibra. E a palavra some.

\noindent Fiz um teste. Escrevi: \textit{\textquotedblleft{}Maria, se você está lendo isto em 1980, acenda uma vela na janela no dia 15 de agosto.\textquotedblright{}} Fechei o diário. Guardei debaixo da cama. No dia 15 de agosto de 1979, liguei para a vizinha. Ela disse: \textquotedblleft{}Sua esposa acendeu uma vela na janela ontem.\textquotedblright{}

\noindent Funciona.

\noindent O diário não é um registro. É uma \textbf{fenda}. E toda fenda tem uma assinatura.

\subsubsection*{\textit{Folha 3 --- 22 de junho de 1979}}

\noindent O padrão apareceu hoje. Passei a noite enviando palavras para datas futuras e medindo o intervalo entre a escrita e o aparecimento.

\noindent O diário não aceita qualquer intervalo. Ele rejeita. A página fica em branco, a tinta escorre, a caneta risca sozinha um risco seco. Funciona apenas em múltiplos de \textbf{62}.

\noindent 62 dias. 124 dias. 186 dias. 248 dias. 310 dias. 372 dias --- e assim por diante.

\noindent Sessenta e dois. O número de anos que Joaquim passou no Colônia. A entidade gravou o número no diário como a assinatura dela. Cada palavra que atravessa o tempo viaja em saltos de 62. É a respiração dela. O ciclo.

\noindent Uma carta destinada a 2026 precisa viajar em saltos de 62 dias ao longo de 47 anos. São 276 saltos. Ou 276 noites de escrita. A página já calculou antes de eu terminar o pensamento --- o número apareceu escrito sozinho na margem direita: \textbf{276}. Ela quer que eu escreva. Ela quer que a carta chegue.

\subsubsection*{Folha 7 --- data ilegível (provavelmente 1985)}

\noindent 2.190 dias se passaram. Escrevo esta nota e a entrego ao diário. Seis anos. Seis ciclos completos.

\noindent O Olho no espelho e vejo a idade chegando. Mas o olho amarelo não envelhece. A cor fica mais forte. Mais densa. Mais \textit{dona}.

\noindent O diário me mostrou uma coisa nova hoje: quando escrevo uma carta destinada a uma data específica, a página 62---sempre a página 62---se ilumina. É ali que a fenda se abre. É ali que a palavra salta.

\noindent A página 62 do diário de Joaquim.

\noindent Ele escreveu aquela página em 1919. Dois anos depois de entrar no Colônia. Ele já tinha descoberto o número. Já sabia que a fome contava em 62.

\noindent Ele escreveu: \textit{\textquotedblleft{}Sessenta e dois é a idade dela. Ela me disse hoje. Nasceu na seca de 1853, na mesma vala, no mesmo cheiro. Sessenta e dois anos depois, eu nasci. Sessenta e dois anos depois, ela me encontrou.\textquotedblright{}}

\noindent A entidade nasceu em 1853. Nasceu da seca. Da fome. Dos corpos empilhados. Ela espera 62 anos. Encontra um hospedeiro. Espera mais 62. Outro hospedeiro.

\noindent O ciclo não é humano. É geológico. É o tempo que a terra leva para esquecer.

\subsubsection*{\textit{Folha 12 --- 2 de março de 2003}}

\noindent O Vinte e quatro anos de ciclos. 151 saltos de 62 dias. A carta está viajando.

\noindent O Às vezes imagino a mão que vai recebê-la. Alguém em 2026. Alguém que vai abrir o arquivo e encontrar um envelope com meu nome. Vai pensar que é uma brincadeira. Depois vai ler a primeira linha.

\noindent E aí já será tarde.

\noindent A carta chegando não é um aviso. É uma entrega. A fenda se abre dos dois lados. Quem lê a carta no futuro também está sendo lido pelo passado.

\noindent Desculpa, doutora. Se é que a senhora existe.

\subsubsection*{Folha 15 --- última folha, 12 de janeiro de 2026}

\noindent O A carta chegou.

\noindent O Eu senti. Hoje, 12 de janeiro de 2026, senti a fenda se fechar. A carta foi recebida. Alguém leu. O ciclo encontrou o próximo.

\noindent O Olho no espelho e vejo um velho que não sou eu. A pele pendurada. O olho direito completamente dourado. A pupila --- já não sei se é minha. Ela se alonga quando o quarto escurece. Fica vertical. Fenda de cabra. De bicho que enxerga no escuro.

\noindent O Eu entreguei o diário hoje. Coloquei num pano preto. Deixei na recepção. A mulher da recepção perguntou meu nome. Eu disse: \textit{\textquotedblleft{}Diga que foi o paciente 1.261.\textquotedblright{}}

\noindent O Ela não entendeu.

\noindent O Mas a pessoa certa vai entender.

\subsubsection*{\textit{Folha 16 --- 12 de janeiro de 2026 (à noite)}}

\noindent O Voltei para casa depois da entrega. Sentei na cama. Olhei minhas mãos. 47 anos. 47 anos vivo contra minha vontade. A carta que escrevi para Maria em 1979 dizia que eu ia proteger os outros. Que eu não ia deixar o diário passar adiante. Eu acreditei nisso por décadas. Acreditei que estava no controle.

\noindent O Não estava. Nunca estive.

\noindent O A fome me manteve vivo. Não por bondade. Não por acaso. Ela precisava de um entregador. Alguém que pudesse andar entre os vivos, que tivesse documento, identidade, aparência de gente. Alguém que chegasse numa recepção de arquivo público em 2026 e não levantasse suspeitas. Um velho de 81 anos, de terno puído, com um pacote embaixo do braço. Invisível. Perfeito.

\noindent O Ela planejou isso desde o dia em que li o diário, em 1979. Desde o dia em que Joaquim escreveu a página 62, em 1919. Desde o dia em que a vala começou a pulsar, em 1853.

\noindent O O ciclo não me perguntou. O ciclo só me usou.

\noindent O Eu pensei que estava protegendo Maria. Pensei que estava quebrando a corrente. Mas a corrente nunca dependeu de mim. Eu era só um elo. E elo não escolhe. Elo segura. Elo passa. Até chegar a vez de outro elo.

\noindent O Maria morreu em 2004. Eu não fui ao enterro. Não porque não quis. Porque não \textit{pude}. No dia do enterro, acordei num banco de praça em outra cidade, com as roupas trocadas e 62 dias de memória perdidos. Ela tinha me levado para longe. Ela não ia deixar ninguém interromper a entrega.

\noindent O Desculpa, Maria. Desculpa, doutora que vai ler isto. Desculpa, próximo da fila.

\noindent O Eu não tive escolha.

\noindent O E vocês também não vão ter.

\noindent \textit{Nota do arquivista: este maço de anotações foi encontrado em fevereiro de 2026, em um envelope pardo lacrado, dentro da mesma caixa que o diário do paciente 1.260. A caligrafia das primeiras folhas é firme, médica. A caligrafia das últimas é trêmula, irregular, com traços que lembram --- mas não idênticos --- à caligrafia de Joaquim Antônio Correia. A data da última anotação (12 de janeiro de 2026) coincide com a data em que o diário foi devolvido ao Arquivo Público Mineiro pela primeira vez.}
\noindent A página 62 do diário original, onde Joaquim escreveu sobre a origem da entidade em 1853, apresenta uma mancha de umidade que cobre parcialmente o texto. A data \textquotedblleft{}1853\textquotedblright{} está escrita três vezes, em caligrafias diferentes --- como se tivesse sido confirmada por mãos diferentes ao longo do tempo. A última confirmação está em letra de médico.

\vspace{0.5em}
\noindent\textit{\textbf{↪ Rotas:} \rota{DOC-01} \rotasep \rota{LUC-02} \rotasep \rota{DOC-27}}
